\documentclass{article} % For LaTeX2e
\usepackage{iclr2027_conference,times}

% Optional math commands from https://github.com/goodfeli/dlbook_notation.
\input{math_commands.tex}

\usepackage{hyperref}
\usepackage{url}
% todonotes needs enough margin space for margin pars.
\setlength{\marginparwidth}{2cm}
\usepackage{todonotes}

% Math packages + a few macros used in the draft
\usepackage{amsmath,amssymb}
\usepackage{xcolor}

% Turn comments on/off:
\newif\ifcomments
\commentstrue      % show comments
%\commentsfalse    % hide comments

% Per-author comment commands (4 authors)
\newcommand{\aditya}[1]{\ifcomments\textcolor{blue}{[\textbf{Aditya}: #1]}\fi}
\newcommand{\claas}[1]{\ifcomments\textcolor{magenta}{[\textbf{Claas}: #1]}\fi}
\newcommand{\siddhant}[1]{\ifcomments\textcolor{teal}{[\textbf{Siddhant}: #1]}\fi}
\newcommand{\amy}[1]{\ifcomments\textcolor{orange}{[\textbf{Amy}: #1]}\fi}

% Dataset / expectation / spaces
\providecommand{\D}{\mathcal{D}}
\providecommand{\E}{\mathbb{E}}
\providecommand{\Sspace}{\mathcal{S}}
\providecommand{\supp}{\operatorname{supp}}
\providecommand{\dd}{\mathrm{d}}


\title{PSMFLows}

% Authors must not appear in the submitted version. They should be hidden
% as long as the \iclrfinalcopy macro remains commented out below.
% Non-anonymous submissions will be rejected without review.

\author{%
Aditya Mohan\thanks{Equal contribution.} \\
Leibniz University Hannover
\And
Claas Voelcker\footnotemark[1] \\
The University of Texas at Austin
\And
Siddhant Agarwal \\
The University of Texas at Austin
\And
Amy Zhang \\
The University of Texas at Austin
}


% The \author macro works with any number of authors. There are two commands
% used to separate the names and addresses of multiple authors: \And and \AND.
%
% Using \And between authors leaves it to \LaTeX{} to determine where to break
% the lines. Using \AN\bibliography{bib/strings}D forces a linebreak at that point. So, if \LaTeX{}
% puts 3 of 4 authors names on the first line, and the last on the second
% line, try using \AND instead of \And before the third author name.

\newcommand{\fix}{\marginpar{FIX}}
\newcommand{\new}{\marginpar{NEW}}

%\iclrfinalcopy % Uncomment for camera-ready version, but NOT for submission.
\begin{document}

\maketitle

\section{Prewriting Form}

1) What problem does this paper solve?

Unsupervised reinforcement learning allows us to train general purpose value functions for a large set (or all) tasks which could be defined in an environment.
Existing methods assume that we want to -- and more importantly, that we can -- train a value or occupancy model for all possible tasks and optimal policies, but that requires a dataset which exhaustively covers the state-action space.
In reality, datasets have limited coverage, and so we need a technique to train a successor basis that is valid for tasks which are well-covered by the dataset.
In this paper, we provide such a technique.

2) Who is the target audience?  What sub-community will be interested in this work?

This paper will be of interest to the community focused on behavior foundation models and unsupervised RL (I believe, Peter, that you don't like the term, but I'm speaking to the audience here ;) ).
Techniques from this community have seen rising interest in robotics domains where large-scale, dense coverage datasets can be easily created in simulation, but it has found limited adoption in data constrained settings such as real world manipulation.
This paper presents an important step towards applicability of successor measure approximations in those scenarios.

3) What are the core technical and conceptual contributions?

Our core technical contribution is twofold: first, we identify how we can define an unconstrained action space for unsupervised RL from a pretrained behavior cloning model.
This allows us to train a PSM style model as if we had access to a dataset that covers all possible actions.
To obtain the necessary deterministic basis policies for PSM training, we fix the noise of the action input. This will give us a large scale set of diverse deterministic policies which remain in the dataset support.
For this, we leverage the fact that we can define the successor measure in the input space of the flow policy, and by construction each latent action will decode to a provided demonstration action.
Secondly, we extend the PSM formalism to this regime (which is relatively simple to do, but necessary).
Finally, we provide empirical evidence that this works well.

4) What are the technical justifications for taking this approach?

Our main goal is to get rid of additional loss regularization terms that are present in other work which need to be tuned to the correct level of pessimism. These are inherited from offline RL.
Our method leverages insights from the DSRL line of work by Andrew Wagenmaker, which has shown to be very robust and reliable in online applications.
Our contributions make these techniques more applicable in offline settings, and show a cool extension to successor measure estimation.

5) What can we now do that wasn’t possible before? What are the benefits of this method?

We can robustly train PSM on limited coverage datasets. Therefore, we can extend PSMs core ideas to domains where getting sufficiently high-fidelity simulation data is difficult, such as real world manipulation.

6) What are some previous methods and why couldn’t they do the same things?

Prior methods either require sim to get very dense datasets, or need to add a bunch of behavior cloning constraints to the learning, which increases the difficulty of hyperparameter tuning and complicates the algorithms.

7) For each contribution / benefit of the method(s), describe the experiment(s) that will showcase this benefit clearly.

Flow inversion: pedagogical experiments that show how a simple flow model (2d can be inverted), larger-scale (OGBench) experiments that show that we can select policies which remain in the data support by sampling from a flow policy, and that we can get diverse deterministic policies by fixing the input noise over multiple steps.

PSM: Show that we can get zero- and few-shot value training on goal reaching tasks in OGBench

8) Why would somebody else use this work or build on it?

Because they want to leverage successor measure approximations e.g. at the VLA scale.

9) What are the limitations of the work, and for each limitation, how will it be addressed (this is also a good place to enumerate any assumptions made, if you haven’t yet)?

Experimental evidence will be limited to ogbench, although we could go for something cooler if we don't make ICLR.
We will likely not get theoretical guarantees from the way we construct the deterministic base policies as we cannot get \emph{all} base policies that are representable (we could, it would just be far more messy).
Flow inversion has some tricky hyperparameters that need to be tuned per task and pretrained policy.

10) What is the single most important point that you want a reader to take away from this paper?

We can do really cool PSM style things by working with a set of pre-trained BC policies, and flows are an excellent way to generate these.


\input{content/abstract}
\section{Introduction}

An agent that quickly adapts to a large variety of tasks in an environment needs to capture a useful representation which encodes how its behavior affects the environment and its own state within it.
To accomplish this, various unsupervised reinforcement learning techniques such as goal-conditioned reinforcement learning, unsupervised skill learning, forward-backward representation learning, and proto successor measure approximations have been proposed.
Recently, \citet{agarwal-rlj25a} have shown that these objectives can be unified as approximations to the successor measure, a generalization of the future state-action occupancy of an agent.

Explicit successor measure approximations such as proto successor measures~\citep{agarwal-arxiv25a} and forward-backward representations~\citep{touati-neurips21a,touati2023does,trinzoni-iclr25a,bagatella-iclr26a} need to be defined with regard to a set of basis policies or tasks.
These prior approaches have explicitly sought to obtain representations capable of representing the occupancy under all possible policies.
However, capturing the behavior of all possible policies from a static dataset is in many cases intractable as the amount of data necessary to accomplish this is simply too large.
Some practical approaches have therefore added ad-hoc regularization to the existing behavior foundation model framework, such as behavior cloning losses. %%AZ.8.29: citation? 

 %%AZ.8.29: Need a link here explaining what's wrong with the above approaches, why should someone prefer the below.
We establish a principled way of extracting a diverse set of well-covered basis policies from a static dataset and obtain a successor measure approximation for this dataset.
For this, we use an expressive flow-based policy to capture the dataset distribution. 
Following \citet{wagenmaker}, we treat the input noise of the flow model as a latent action, and define deterministic basis policies in this latent space.
We can then treat this noise space as an unconstrained action space and any policy defined over it will stay close to the set of policies well-covered by the data.

 %%AZ.8.29: the contribution of this paragraph is not coming through. Is the point that we are trying to augment the latent action dataset? What is the benefit of doing this?
However, while flow policies allow us to define unconstrained basis policies, the data itself contains environment actions, not latent actions.
Multiple latent actions can lead to the same environment action, and identifying these allows us to augment the dataset with all valid latent transitions.
We therefore present an algorithm to identify a distribution over latent noise actions which are decoded by the flow policy to similar environment action.
This technique allows us to augment the training dataset for the proto successor measures, but we believe it is also of general interest to the community for many other applications.
To construct the noise space sets, we leverage the local invertibility of the flow model.
We then compute a proposal distribution in the latent space and treat the set identification as a posterior inference problem.

The resulting algorithm for successor measure approximation using inverted flow models and latent actions is called Proto Successor Measure Flows, a principled method for computing successor approximations over limited coverage datasets.
We implement and validate our technique on the OGBench tasks, a goal-conditioned environment and dataset that naturally allows for testing multi-task goal-reaching capabilities of behavior foundation models.
We show ... \todo[inline]{add results here}.
\input{content/related-work}
\section{Preliminaries}

\todo[inline]{add markov boilerplate}

\subsection{Markov Decision Processes and Successor Measures}

For a policy \(\pi\), define the discounted successor measure
\[
M^\pi(s,\mathrm{d}x) = \E_\pi \left[ \sum_{t=0}^\infty \gamma^t \, \1_{\{s_t\in \mathrm{d}x\}} \;\middle|\; s_0=s \right].
\]
For any reward \(r:\Sspace\to\R\),
\[ V_r^\pi(s) = \int r(x)\, M^\pi(s,\mathrm{d}x). \]

Let \( \D={(s_i,a_i,s_i')}_{i=1}^N \) be a fixed offline dataset collected by an unknown behavior policy \(\mu_{\D}(a\mid s)\).  
We want to learn SMs on \(\D\).

The offline difficulty is that Bellman backups and policy extraction require action queries such as \( \max_{a'} Q(s',a'), \) which may select
\[
a' \notin \supp \mu_{\D}(\cdot\mid s').
\]
In SMs, this causes the representation to be trained or queried outside the dataset-supported action manifold.

Therefore, we require a mechanism to learn successor measures $M^\pi$ for all $\pi \in \Pi_c$, where $\Pi_c$ is the set of all policies which are well covered by the dataset.
For a policy to be ``well-covered'' we require that \todo[inline]{formalize that a state action pair visited by the policy has a high chance of being captured in the dataset}


\subsection{The Proto Successor Measure Framework}

Furthermore, setups like FB entangle "task" indicator and "policy" indicator, implicitly constraining policies to the set of all optimal policies for all tasks.

PSMs represent the space of possible SMs using policy-independent basis measures:

\[
    M^\pi \approx b + P^\top w^\pi ,
\]


where $b$ and $P$ are sheared across policies, and \(w^\pi\) are policy-specific coordinates.

For a reward $r$, policy extraction becomes: 
\[
    w_r^\star \in \argmax_w \langle b+P^\top w,\ r\rangle,
\]
subject t the constraint of \(b+P^\top w\).
Operationally, PSM uses a hashing mechanism \( z \mapsto \pi_z \).

The central problem, then, is how to construct a useful policy class \(\{\pi_z \}\) from limited-support offline data.

The trick is to obtain a space in which an easy to sample distribution over $z$ directly indexes into the constrained policy space.

\subsection{Behavior Cloning with Flow Policies}

Train a conditional behavior flow on \(\D\). 
For each dataset action $a$, sample

\[
    u_0  \sim p_0. \qquad u_1 = a, \qquad (s,a) \sim \D
\]

Here $u_0$ is the initial flow noise and $u_1$ is the target action.
For $t \in [0,1]$, define
\[
    u_t = (1-t)u_0 + tu_1
\]
Th target velocity along this path is
\[
    \dot u_t = u_1 - u_0 = a - u_0 
\]

We train a conditional vector field 
\[
    v_\theta:\Sspace \times [0,1] \times \R^{d_u} \to \R^{d_u}, \qquad (s,t,u_t)\mapsto v_\theta(s,t,u_t),
\]
with the BC flow objective
\[
    \mathcal L_{\mathrm{flow}} (\theta) = \E_{(s,a) \sim \D,\ u_0\sim p_0,\ t\sim \mathrm{Unif}([0,1])}
    \left[ \left\| v_\theta(s,t,u_t) - (u_1-u_0) \right\|_2^2 \right].
\]

At inference, givem a state $s$ and latent $u \sim p_0$, the flow is obtained by solving
\[
    \frac{\dd u_t}{\dd t} = v_\theta(s,t,u_t), \qquad u_{t=0}=u .
\]
\section{The space of policies indexed by $u$}

To define a set of restricted policy spaces, we will now define policies over the set of initial noise variables $u_0$ sampled from the flow prior as $\Pi = \{G_\theta(s,u_0)|u_0 \sim p_0\}$.

This set is not equal to the set of all well-covered policies, importantly any policy that resamples $u_0$ at every step is also a valid flow policy under the data.
However, representing this would require a significantly larger set of latent variables.

We can now train a PSM approximation as $\phi(s,u_0,u'_0)^T\psi(s^+)$. Note that $\phi$ takes two $u$ as an input: one as the current state action, and one that is the index of the future rollout.
Obtaining a meaningful $u_0$ from dataset samples $(s,a,s')$ requires us to invert the flow policy (which seems feasible given the recent literature \url{https://arxiv.org/pdf/2605.10821})\footnote{We probably want to expand this, because this just generates one noise sample for each action, but we want all noise samples which give us a similar action, so we need to add some noise, and maybe a forward correction pass.}.

The interpretation of $\phi(s,u_0,u_0')$ is now as the feature for a taking action $a=G_\theta(s,u_0)$ and then following $\pi = G_\theta(s,u_0')$ afterwards.

\section{Policy inference}

Given a task expressed via weights $w$, we can now run a constrained Q learning procedure to obtain a Q function for \textbf{any} well-covered mixture policy without leaving the $z$ space in which all of our actions remain in distribution.
We need to solve for the following recursive function on the dataset (again, generating $u_0$ from provided actions)
\begin{align}
    Q^*(s,u_0,w) = \max_{u'_0} \left(\phi(s,u_0,u_0') - \gamma \phi(s',u_0',u_0')\right)^\top w + \gamma Q^*(s',u_0',w)
\end{align}

In principle $\left(\phi(s,u_0,u_0') - \gamma \phi(s',u_0',u_0')\right)^\top w$ should actually be independent of $u_0'$ if I'm not completely mistaken because it is simply the one step occupancy / one step reward? This would simplify the problem to a standard Q learning problem:

\begin{align}
    Q^*(s,u_0,w) = \left(\phi(s,u_0,u_0) - \gamma \phi(s',u_0,u_0)\right)^\top w + \max_{u'_0} \gamma Q^*(s',u_0',w).
\end{align}

We can solve for this Q function and essentially obtain a latent noise DSRL style critic for free. Now finally we need to sample initial noise to maximize $Q^*$ for the final policy.
We can even use Haoran's technique here in the z space!

Alternatively, we can train Q on intermediate noise samples as well $u_t$ for $t\in[0,1]$, and use the Q function instead as a steering approach. 

\input{content/experiments}

\subsection*{AI use statement}

(This section is \textbf{required} and does not count toward the page limit.)

In this work, we used generative AI tools for [tasks with required disclosure].
We have not used generative AI tools for [other tasks with required disclosure],
and [the rest of the required disclosure tasks] are not applicable to this work.
Additionally, we used generative AI tools for [tasks with recommended
disclosure]. We have reviewed all AI-assisted work. [Elaborate. For example, “we
checked LLM-generated research ideas for potential plagiarism through a manual
literature survey”, “LLM-generated code was verified and tested for correctness
by 2 authors”, etc.]. We take responsibility for the final content of this work,
including text, claims or artifacts produced with the aid of generative AI.

See the ICLR 2027 AI Policy for Authors for more details. This statement should
not be more than 1 page.

\subsection*{Ethics statement}

(This section is \textbf{recommended} and does not count toward the page limit.)

If authors feel that their paper submission raises questions regarding the Code
of Ethics, they are encouraged to include a paragraph of Ethics Statement (at
the end of the main text before references) to address potential concerns where
appropriate. Topics include, but are not limited to, studies that involve human
subjects, practices to data set releases, potentially harmful insights,
methodologies and applications, potential conflicts of interest and sponsorship,
discrimination/bias/fairness concerns, privacy and security issues, legal
compliance, and research integrity issues (e.g., IRB, documentation, research
ethics). This statement should not be more than 1 page.

\subsection*{Reproducibility statement}

(This section is \textbf{recommended} and does not count toward the page limit.)

It is important that the work published in ICLR is reproducible. Authors are
strongly encouraged to include a paragraph-long Reproducibility Statement at the
end of the main text (before references) to discuss the efforts that have been
made to ensure reproducibility. This paragraph should not itself describe
details needed for reproducing the results, but rather reference the parts of
the main paper, appendix, and supplemental materials that will help with
reproducibility. For example, for novel models or algorithms, a link to an
anonymous downloadable source code can be submitted as supplementary materials;
for theoretical results, clear explanations of any assumptions and a complete
proof of the claims can be included in the appendix; for any datasets used in
the experiments, a complete description of the data processing steps can be
provided in the supplementary materials. Each of the above are examples of
things that can be referenced in the reproducibility statement.


\subsubsection*{Acknowledgments}
Use unnumbered third level headings for the acknowledgments. All
acknowledgments, including those to funding agencies, go at the end of the paper.


\bibliographystyle{iclr2027_conference}
\bibliography{bib/local,bib/lib,bib/proc,bib/strings,iclr2027_conference}

\appendix
\section{Appendix}
You may include other additional sections here.



\end{document}

